\section{An arbitrary-order image hierarchy}\label{sec:arbitrary}

The identity
\begin{equation}\label{eq:sinh-integral}
 \frac{a}{\theta}\sinh\frac{a\theta}{t}
 =\frac{a^2}{t}\int_0^1\cosh\frac{sa\theta}{t}\,\dd s
\end{equation}
removes the apparent singularity in \eqref{eq:q} before differentiation.

Fix $M\in\N_0$, $0<\rho<\pi$, and $t_0>0$.  For $0\le m\le M$ define
\begin{equation}\label{eq:Qm}
 Q_m(t_0,\rho)=2\sum_{k\ge1}
 e^{-a_k^2/(2t_0)+a_k\rho/t_0}
 \left(\frac{a_k}{t_0}\right)^m
 \left(1+\frac{a_k^2}{t_0(m+1)}\right).
\end{equation}
The series is explicit and computable.  If it is stopped after $K$ terms,
the omitted tail is at most the first omitted term divided by $1-r_{m,K+1}$,
where
\begin{equation}\label{eq:ratio}
r_{m,N}=\left(\frac{N+1}{N}\right)^{m+2}
\exp\left[-\frac{2\pi^2(2N+1)-2\pi\rho}{t_0}\right].
\end{equation}

\begin{theorem}[Uniform arbitrary-order image enclosure]\label{thm:image-hierarchy}
Suppose $0<\rho<\pi$ and
\begin{equation}\label{eq:time-condition}
 t_0\le \frac{(2\pi)^2/2-2\pi\rho}{M+1}.
\end{equation}
Then, for every $0<t\le t_0$, $0\le\theta\le\rho$, and
$0\le m\le M$,
\begin{equation}\label{eq:qbound}
 \abs{q_t^{(m)}(\theta)}\le Q_m(t_0,\rho).
\end{equation}
For each fixed $M,t_0,\rho$, the right-hand side has a rigorously bounded geometric
tail of the form \eqref{eq:ratio}.
\end{theorem}

\begin{proof}
Differentiate \eqref{eq:sinh-integral} under the integral.  Each hyperbolic
derivative is bounded by $e^{a\rho/t}$, while the integral contributes
$\int_0^1s^m\dd s=(m+1)^{-1}$.  This gives the $k$th summand in
\eqref{eq:Qm}.  Writing its $t$-dependence as
$e^{-A_k/t}t^{-m}(1+c_k/t)$ shows that it increases with $t$ whenever
$t\le A_k/(m+1)$.  The smallest $A_k=a_k^2/2-a_k\rho$ occurs at $k=1$;
\eqref{eq:time-condition} therefore makes $t=t_0$ the uniform worst case.
Finally, the quotient of successive summands is bounded by
\eqref{eq:ratio}; details appear in Appendix~\ref{app:image-tails}.
\end{proof}

The theorem is arbitrary-order in the standard effective sense: for every
requested finite $M$ it gives a closed computation.  It is not a claim that a
single fixed $t_0$ satisfies \eqref{eq:time-condition} for all $M$ at once.
